\section{DeepSeekMath-V2: que el verifier también acepte ser verificado}

\subsection{El outcome reward no puede garantizar el proceso de la demostración}

\videofigure{figures/outcome-reward-models.jpg}{El outcome reward solo revisa la respuesta final; no puede determinar si la demostración intermedia tiene pasos omitidos, razonamiento circular o un acierto por suerte.}{00:15:02--00:16:52}

\videofigure{figures/proof-verification-problem.jpg}{Un LLM puede generar demostraciones con forma fluida pero matemáticamente inválidas, y sin una solución de referencia, incluso un judge ordinario tiene dificultad para ubicar el problema de manera estable.}{00:16:52--00:18:18}

Una demostración matemática no es un simple emparejamiento de respuesta corta. Un theorem proof puede tener una conclusión correcta, pero algún paso no se deduce de las premisas; si solo se da un reward final, esa flawed trajectory se tratará como muestra positiva. El process reward model, por su parte, depende de etiquetas humanas paso a paso, que son costosas.

\begin{warningbox}{El verifier puede tener "puntuación correcta, razón equivocada"}
Aunque un LLM judge dé una puntuación global aproximadamente correcta, puede inventar errores que no existen. Si esas critique se usan para entrenar al generator, el sistema termina escribiendo de vuelta en el modelo las alucinaciones del verifier.
\end{warningbox}

\subsection{Generator--Verifier--Meta-Verifier}

\videofigure{figures/deepseekmath-v2-architecture.jpg}{DeepSeekMath-V2 construye un ciclo cerrado de colaboración entre generator, verifier y meta-verifier.}{00:18:18--00:19:34}

La estructura puede formalizarse como:

$$
y\sim G_\theta(x),\qquad
(e,s)=V_\phi(x,y),\qquad
m=M_\psi(x,y,e,s),
$$

donde el generator produce la demostración $y$, el verifier da un análisis del problema $e$ y una puntuación $s$, y el meta-verifier determina si el problema realmente existe y si la puntuación se deduce de ese análisis.

\videofigure{figures/meta-verification.jpg}{El meta-verifier hace una segunda revisión del análisis del verifier, reduce los fabricated issue y aumenta la confiabilidad de la explicación.}{00:19:34--00:20:18}

\begin{importantbox}{La meta-verification revisa la cadena de evidencia de la evaluación}
No es simplemente votar una vez más, sino verificar si cada issue de la critique puede ubicarse en la demostración original, y si el score final es consistente con esos issue.
\end{importantbox}

\subsection{Optimización iterativa y resultados}

\videofigure{figures/evaluation-methodology.jpg}{El sistema combina GRPO, generación--verificación iterativa, y evalúa la capacidad de autoverificación con varios tipos de benchmark de demostraciones matemáticas.}{00:20:18--00:20:52}

\videofigure{figures/proof-results.jpg}{En IMO Shortlist 2024, tanto pass@1 como best@32 aumentan de manera sostenida con más rondas de iteración de self-verification.}{00:20:52--00:21:22}

\videofigure{figures/self-verification-improvement.jpg}{Las demostraciones con el verifier score más alto se seleccionan para la siguiente ronda de entrenamiento, formando un ascenso conjunto entre el generator y el verifier.}{00:21:22--00:22:22}

Si se elige, de entre $K$ demostraciones, la que tenga la puntuación de verifier más alta:

$$
y^*=\arg\max_{y_k,\,k\le K}V_\phi(x,y_k),
$$

entonces la premisa es que el ranking de $V_\phi$ sea consistente con la calidad real de la demostración. Precisamente el papel de la meta-verification es mejorar ese alignment, para que best@$K$ no se limite a elegir la demostración que "mejor le agrada al judge".

\begin{warningbox}{Más capas de judge no traen la verdad de manera automática}
El generator, el verifier y el meta-verifier pueden compartir datos de entrenamiento y sesgos, produciendo correlated error. Aun así se necesitan revisiones formales, modelos independientes, auditorías de expertos o adversarial test para romper esa complicidad.
\end{warningbox}

\subsection{Resumen de la sección}

DeepSeekMath-V2 transforma al verifier de un árbitro incuestionable a un objeto auditable. Muestra una vía de verificación automática del proceso cuando no hay una demostración de referencia, pero sigue dependiendo de la estructura lógica relativamente clara del dominio matemático.
